% =====================================================================
%  theory/interface_conditions.tex
%  Phase 1, Blueprint V5.0, Task D
%
%  Derivation of the generalized imperfect-contact conditions, their
%  dissipation forms, passivity, and the INDEPENDENT EQUATION COUNT for
%  each configuration.
%
%  SOURCE-ONLY: no TeX engine is installed. Verified as text; the symbolic
%  claims are checked by code/foundation_checks.py (checks F1, I2).
% =====================================================================
\documentclass[11pt]{article}
\usepackage{amsmath,amssymb,amsthm,bm}
\usepackage{geometry}
\usepackage{booktabs}
\geometry{margin=1in}
\newcommand{\s}{\sigma}
\newcommand{\bth}{\vartheta}
\newcommand{\D}{\bm{D}}
\newcommand{\E}{\bm{E}}
\newcommand{\bu}{\bm{u}}
\newcommand{\bv}{\bm{v}}
\newcommand{\bq}{\bm{q}}
\newcommand{\bt}{\bm{t}}
\newcommand{\jj}{[\![}          % jump open
\newcommand{\kk}{]\!]}          % jump close
\newcommand{\jmp}[1]{[\![#1]\!]}
\theoremstyle{plain}
\newtheorem{proposition}{Proposition}
\newtheorem{remark}{Remark}
\newtheorem{definition}{Definition}

\title{Generalized imperfect interface conditions, dissipation and\\
equation counting\\
\large Theory note 2 of 4 --- Phase 1, Blueprint V5.0}
\author{Derived for the Piezothermoelastic Anti-Reflection programme}
\date{4 November 2026}

\begin{document}
\maketitle

\begin{abstract}
This note derives the contact laws of Blueprint V5.0 \S4 together with the
interface power balance they imply, and it counts the independent interface
relations for each configuration. Three things are established rather than
assumed: (i) the mechanical, thermal and (where applicable) electrical contact
dissipations are non-negative forms that follow from the contact laws plus a
passivity requirement, with units checked; (ii) the $\mathrm{h}\to 0$ composition
of the interphase is a face-compliance map, not a welded contact, and two
compliant faces do not give zero jump; (iii) the number of independent relations
equals the number of unknown amplitudes after gauge removal, and this count is
$8$ for the two-half-space baseline and $16$ for the interphase configuration.
The passivity claims are verified symbolically/numerically by checks F1 and I2 of
\texttt{code/foundation\_checks.py}. \textbf{No reflection or transmission
coefficient is computed in this note.}
\end{abstract}

\section{Setting and notation}
\label{sec:setting}

Two half-spaces meet at a face. Medium 1 occupies $z<0$, medium 2 occupies $z>0$
(baseline) or $z>\mathrm{h}$ (interphase configuration, with $0<z<\mathrm{h}$
occupied by a third medium). On a plane $z=\text{const}$ define the state and the
generalized flux
\begin{equation}
  \bm{q} = (u_x,\;u_z,\;\varphi,\;\bth)^{\mathsf T},
  \qquad
  \bm{F} = (\s_{xz},\;\s_{zz},\;D_z,\;q_z)^{\mathsf T},
  \label{eq:state-flux}
\end{equation}
in the notation of Blueprint \S6. The mechanical traction vector is
$\bt=(\s_{xz},\s_{zz})^{\mathsf T}$, the velocity is $\bv=\dot{\bu}$, and $q_z$
is the heat flux in $+z$.

The jump is defined with the ordering ``2 minus 1'',
\begin{equation}
  \jmp{f} \equiv f^{(2)} - f^{(1)} .
  \label{eq:jump}
\end{equation}
This ordering is fixed once here; the sign of every contact law below follows
from it and must not be flipped independently.

\section{Mechanical contact}
\label{sec:mech}

\subsection{The law}

\begin{definition}[Passive diagonal spring--dashpot contact]
\label{def:spring}
Tractions are continuous and the velocity jump is carried by a diagonal
spring--dashpot element:
\begin{equation}
  \bt^{(1)} = \bt^{(2)} = \bt_c,
  \qquad
  \bt_c = K\,\jmp{\bu} + C\,\jmp{\bv},
  \qquad K = \mathrm{diag}(K_x,K_z)\succeq 0,\quad
  C = \mathrm{diag}(C_x,C_z)\succeq 0 .
  \label{eq:spring}
\end{equation}
\end{definition}

This is Blueprint Eq.~(6). The positivity constraints $K\succeq0$, $C\succeq0$
are the definition of a \emph{passive} contact: no energy is supplied by the
interface. Note that \eqref{eq:spring} is a statement about the \emph{traction},
so the element carries the same traction on both of its sides --- this is what
makes the dissipation calculation below come out non-negative.

\subsection{Derivation of the interface dissipation}
\label{sec:mech-diss}

\begin{proposition}[Mechanical contact dissipation]
\label{prop:mech-diss}
For the contact \eqref{eq:spring} the mechanical power absorbed at the interface,
defined as the difference between the power flowing into the contact and the
power flowing out of it, is
\begin{equation}
  \mathcal{D}_{\mathrm{mech}}
  = \tfrac12\,\jmp{\bv}^{\mathsf T} C \jmp{\bv}
  \;+\; \tfrac12\,\jmp{\bu}^{\mathsf T} K\,(-i\omega)\,\jmp{\bu}^{*}
  \quad\text{(time-harmonic, }e^{-i\omega t}\text{)},
  \label{eq:mech-diss}
\end{equation}
of which the first term is the genuinely dissipated part and the second is
purely reactive.
\end{proposition}

\begin{proof}
The power per unit area entering the contact from medium 1 is
$-\bt^{(1)}\cdot\bv^{(1)}$ (the minus sign is because the outward normal of
medium 1 at the face is $+z$ while the traction is transmitted in $+z$). The power
leaving into medium 2 is $-\bt^{(2)}\cdot\bv^{(2)}$. The absorbed power is their
difference,
\begin{equation}
  \mathcal{D} = -\bt^{(1)}\cdot\bv^{(1)} + \bt^{(2)}\cdot\bv^{(2)}
  = -\bt_c\cdot\big(\bv^{(2)}-\bv^{(1)}\big)
  = -\bt_c\cdot\jmp{\bv},
  \label{eq:mech-diss-step1}
\end{equation}
using $\bt^{(1)}=\bt^{(2)}=\bt_c$ from \eqref{eq:spring}. Substituting
\eqref{eq:spring},
\begin{equation}
  \mathcal{D} = -\big(K\jmp{\bu} + C\jmp{\bv}\big)\cdot\jmp{\bv}
  = -\jmp{\bu}^{\mathsf T}K\jmp{\bv} - \jmp{\bv}^{\mathsf T}C\jmp{\bv}.
  \label{eq:mech-diss-step2}
\end{equation}
With $e^{-i\omega t}$, $\bv = -i\omega\bu$, so $\jmp{\bv} = -i\omega\jmp{\bu}$ and
$\jmp{\bu}^{\mathsf T}K\jmp{\bv}
 = -i\omega\,\jmp{\bu}^{\mathsf T}K\jmp{\bu}$. The cycle average of the first
term of \eqref{eq:mech-diss-step2} is therefore
$-\operatorname{Re}\langle -i\omega\,\jmp{\bu}^{\mathsf T}K\jmp{\bu}\rangle$,
which is purely reactive (imaginary in the phasor sense, vanishing on average)
because $K$ is real and symmetric. The second term is
$\jmp{\bv}^{\mathsf T}C\jmp{\bv} = \omega^2\jmp{\bu}^{\mathsf T}C\jmp{\bu}$,
which is real, positive semi-definite because $C\succeq0$, and is the true
dissipation. Time-averaging the harmonic product gives the factor $\tfrac12$,
yielding \eqref{eq:mech-diss}.
\end{proof}

\begin{proposition}[Non-negativity]
\label{prop:mech-nonneg}
For $C\succeq0$ the dissipated part is non-negative for every velocity jump:
\begin{equation}
  \tfrac12\,\jmp{\bv}^{\mathsf T} C \jmp{\bv} \ge 0,
  \qquad \text{units } \mathrm{W\,m^{-2}} .
  \label{eq:mech-nonneg}
\end{equation}
\end{proposition}
\begin{proof}
$C$ is real symmetric positive semi-definite, so $\jmp{\bv}^{\mathsf T}C\jmp{\bv}\ge0$.
The unit check: $[\jmp{\bv}]=\mathrm{m\,s^{-1}}$ and (for the dashpot)
$[C]=\mathrm{Pa\,s\,m^{-1}}=\mathrm{N\,s\,m^{-3}}$, so
$[C]\,[\jmp{\bv}]^2 = \mathrm{N\,s\,m^{-3}}\cdot\mathrm{m^2\,s^{-2}}
= \mathrm{N\,m^{-1}\,s^{-1}} = \mathrm{W\,m^{-2}}$. Correct.
\end{proof}

\textsc{Tested}: check F1 draws $5\times10^{4}$ random positive semi-definite
$C$ and random velocity jumps and finds a minimum of exactly $0$ (the samples
include the degenerate case). No negative value is found.

\begin{remark}[The traction sign is not free]
\label{rem:traction-sign}
Blueprint \S5.2 uses the force-like variable $f=-\s_{zz}$, chosen so that a
right-going wave obeys $f=Zv$ with $Z>0$. Equation~\eqref{eq:spring} above is
written in the traction $\bt$ with the \emph{opposite} sign. Both are consistent,
but a calculation that mixes them silently reverses the sign of the reflection
coefficient and therefore of the Voigt mapping $Y_d = 1/(C+iK/\omega)$. Blueprint
\S13 item 1 asks for this to be checked independently; the result is that
Proposition~1 of the blueprint is correct \emph{in its own convention} $(f=-\s_{zz})$
and that the mapping to \eqref{eq:spring} carries one sign flip. This is recorded
as a convention caution, not an error in the blueprint.
\end{remark}

\section{Electrical contact}
\label{sec:elec}

Two physically distinct closures are available, and \textbf{exactly one must be
chosen}: imposing both an arbitrary free surface charge \emph{and} an independent
potential discontinuity over-constrains the problem, which is the trap identified
in Blueprint \S13 item 2.

\subsection{Closure E-A: continuous potential (short-circuit / no interfacial charge storage)}
\begin{equation}
  \jmp{\varphi} = 0, \qquad \jmp{D_z} = 0 .
  \label{eq:elec-EA}
\end{equation}
Two relations.

\subsection{Closure E-B: capacitive gap with no free charge}
\begin{definition}[Passive capacitive interface]
\label{def:cap}
No free charge is present, and the normal electric displacement is carried by a
lumped capacitance per unit area $C_s\ge0$:
\begin{equation}
  \jmp{D_z} = 0, \qquad D_z = C_s\,\jmp{\varphi}, \qquad C_s \ge 0 .
  \label{eq:elec-EB}
\end{equation}
\end{definition}
Two relations. The energy stored per unit area is
$\tfrac12 C_s \jmp{\varphi}^2$, with $[C_s]=\mathrm{F\,m^{-2}}$ and
$\tfrac12\cdot\mathrm{F\,m^{-2}}\cdot\mathrm{V^2}
= \mathrm{J\,m^{-2}}$. Correct.

\begin{proposition}[The capacitive gap is reactive, not dissipative]
\label{prop:cap-reactive}
For the closure \eqref{eq:elec-EB} the cycle-averaged electrical power absorbed
at the interface is exactly zero.
\end{proposition}
\begin{proof}
With no free charge and no conductive path across the gap, the interfacial
current is $i\omega C_s\jmp{\varphi}$ and the absorbed power is
$\tfrac12\operatorname{Re}\big[(i\omega C_s\jmp{\varphi})^{*}\jmp{\varphi}\big]
= \tfrac12\operatorname{Re}(-i\omega C_s)\lvert\jmp{\varphi}\rvert^2 = 0$
because $C_s$ is real. The interface behaves as a pure capacitance: it stores and
returns energy each cycle but never dissipates.
\end{proof}
\textsc{Tested} (check F1 evaluates the expression and obtains exactly $0$).

\begin{remark}[Why this matters for the suppression question]
\label{rem:reactive}
A purely reactive electrical contact cannot \emph{absorb} energy, so it cannot by
itself deliver the dissipative matching of Proposition~1(b) of the blueprint. It
can only shift the modal impedances. Any claim that an electrical contact law
``absorbs'' reflection must therefore identify a conductive or leakage path;
if none is declared, the electrical channel is reactive and the suppression must
come from the mechanical or thermal channel.
\end{remark}

\subsection{If a conductive path is intended}
A leakage conductance $G_s\ge0$ in parallel with the capacitance replaces
$i\omega C_s \to G_s + i\omega C_s$ in \eqref{eq:elec-EB}, and the absorbed power
becomes $\tfrac12 G_s\lvert\jmp{\varphi}\rvert^2 \ge 0$, with
$[G_s]=\mathrm{S\,m^{-2}}$. This is admissible and passive, but it must be
\emph{declared}. It is \textsc{Unresolved} whether any physical interface of
interest here has such a path; Phase 1 does not assert one.

\section{Thermal contact}
\label{sec:therm}

\subsection{The law}

\begin{definition}[Passive thermal contact resistance]
\label{def:thermal}
Heat fluxes are continuous and the temperature jump drives the flux:
\begin{equation}
  q_z^{(1)} = q_z^{(2)} = q_c, \qquad
  q_c = H_T\big(\bth^{(1)} - \bth^{(2)}\big) = -H_T\,\jmp{\bth},
  \qquad H_T \ge 0 .
  \label{eq:thermal}
\end{equation}
\end{definition}
This is Blueprint Eq.~(7). The sign is fixed as follows: if medium 1 is hotter
($\bth^{(1)}>\bth^{(2)}$, so $\jmp{\bth}<0$), heat must flow from 1 to 2, i.e. in
$+z$, so $q_c>0$. Equation~\eqref{eq:thermal} gives
$q_c = -H_T\,\jmp{\bth} > 0$. Correct.

\subsection{Dissipation}

\begin{proposition}[Thermal contact dissipation]
\label{prop:therm-diss}
The entropy production rate at the interface is
\begin{equation}
  \dot S_{\mathrm{int}} = q_c\left(\frac{1}{T^{(2)}} - \frac{1}{T^{(1)}}\right)
  \;\simeq\; \frac{q_c\,\jmp{\bth}}{T_0^{2}}
  \;\overset{\eqref{eq:thermal}}{=}\;
  \frac{H_T\,\lvert\jmp{\bth}\rvert^{2}}{T_0^{2}} \ \ge 0 ,
  \label{eq:therm-entropy}
\end{equation}
so that the corresponding dissipation per unit area, $T_0\dot S_{\mathrm{int}}$,
is
\begin{equation}
  \boxed{\ \mathcal{D}_{\mathrm{th}} = \frac{H_T\,\lvert\jmp{\bth}\rvert^{2}}{T_0}
  \ \ge 0\ }\qquad \text{units } \mathrm{W\,m^{-2}}.
  \label{eq:therm-diss}
\end{equation}
\end{proposition}

\begin{proof}
Entropy is produced whenever heat crosses a finite temperature drop. The flux
$q_c$ crosses from $T^{(1)}$ to $T^{(2)}$, giving
$\dot S = q_c(1/T^{(2)}-1/T^{(1)})$. Linearising about $T_0$,
$1/T^{(2)}-1/T^{(1)} \simeq (T^{(1)}-T^{(2)})/T_0^2 = -\jmp{\bth}/T_0^2$, so
$\dot S \simeq -q_c\jmp{\bth}/T_0^2$, which with \eqref{eq:thermal} gives
$H_T\lvert\jmp{\bth}\rvert^2/T_0^2 \ge 0$. Multiplying by $T_0$ gives
\eqref{eq:therm-diss}. The unit check:
$[H_T]=\mathrm{W\,m^{-2}K^{-1}}$, so
$[H_T][\bth]^2/[T_0] = \mathrm{W\,m^{-2}K^{-1}}\cdot\mathrm{K^2}/\mathrm{K}
=\mathrm{W\,m^{-2}}$. Correct.
\end{proof}

\begin{remark}[On the absolute-temperature factor]
\label{rem:absolute-T}
Blueprint \S4 offers $\tfrac12 H_T\lvert\jmp{\bth}\rvert^2/T_0$ as a
\emph{candidate} and \S13 item 3 requires it to be derived. The derivation above
gives \eqref{eq:therm-diss} \emph{without} the factor $\tfrac12$: the missing
half is because the derivation uses time-averaged squares rather than peak
squares, whereas the mechanical form \eqref{eq:mech-nonneg} is written with the
conventionally included $\tfrac12$. \textbf{Adopting one convention for all three
channels is required before any numerical energy balance}; mixing them produces a
50~\% error in the thermal dissipation. This is flagged as an unresolved
convention item (U-5) rather than quietly resolved, because the correct choice
depends on how $R,T,A$ are normalised in the Phase-2 solver.
\end{remark}

\textsc{Tested}: check F1 samples $2\times10^{4}$ random $(H_T,\jmp{\bth})$ and
finds a minimum of exactly $0$.

\section{The $h\to 0$ reduction}
\label{sec:h0}

Blueprint \S6 requires the interphase composition to be \emph{derived}, with the
expected form marked as a hypothesis. The derivation is as follows.

Write the state $\bm{s}=(\bm{q},\bm{F})$ of \eqref{eq:state-flux} and suppose the
interphase admits a linear evolution
\begin{equation}
  \frac{d\bm{s}}{dz} = \bm{A}(\omega,\kappa)\,\bm{s}, \qquad 0<z<\mathrm{h}.
  \label{eq:state-evolution}
\end{equation}
Each face imposes $4$ relations. Write them as maps
$\bm{P}_i\bm{s} = 0$ ($i=1$ at $z=0$, $i=2$ at $z=\mathrm{h}$). The composite
map from the state just below the lower face to the state just above the upper
face is
\begin{equation}
  \bm{s}(\mathrm{h}^{+}) = \bm{P}_2\,e^{\mathrm{h}\bm{A}}\,\bm{P}_1\,\bm{s}(0^{-}),
  \label{eq:composite}
\end{equation}
which is Blueprint \S6's composition.

\begin{proposition}[First-order reduction]
\label{prop:h0}
As $\mathrm{h}\to0$ the composition \eqref{eq:composite} becomes
\begin{equation}
  \bm{F}^{(2)}-\bm{F}^{(1)} = -\Big[\bm{Y}_1 + \bm{Y}_2 + \mathrm{h}\,\bm{Y}_{\ell}\Big]\,
  \big(\bm{F}^{(1)}+\bm{F}^{(2)}\big)/2 \;+\; O(\mathrm{h}^2),
  \label{eq:h0}
\end{equation}
where $\bm{Y}_1,\bm{Y}_2$ are the face compliances and
$\bm{Y}_\ell = \bm{A}_{Fq}$ the layer's own flux-to-state coupling at zero
frequency. In words: \textbf{the jumps add} and the layer contributes an
$O(\mathrm{h})$ correction.
\end{proposition}

\begin{proof}[Proof sketch, with the status of each step stated]
Expand $e^{\mathrm{h}\bm{A}} = \bm{I} + \mathrm{h}\bm{A} + O(\mathrm{h}^2)$ in
\eqref{eq:composite}, substitute the face laws linearised about the mean state,
and collect terms of equal order in $\mathrm{h}$. The $O(1)$ term is the sum of
the two face compliance maps, which is \eqref{eq:h0} with $\mathrm{h}=0$; the
$O(\mathrm{h})$ term is the layer coupling. \textsc{Derived} to first order.
The second-order term and the question of whether \eqref{eq:h0} can be written
with a single effective compliance for general (non-diagonal) face laws are
\textsc{Unresolved} (U-6); the blueprint's own wording marks this as a hypothesis,
and the hypothesis is confirmed only in the diagonal, normal-incidence case.
\end{proof}

\begin{remark}[Two compliant faces are not a welded contact]
\label{rem:not-welded}
Equation~\eqref{eq:h0} shows that the jump does \emph{not} vanish as
$\mathrm{h}\to0$ unless both face compliances vanish. The limit
$\mathrm{h}\to0$ of the interphase configuration is therefore the
\emph{imperfect-contact} two-half-space baseline, not the welded one. Blueprint
\S6 states this as an expectation; here it is derived to first order. Operationally
this means the interphase configuration must be compared against \emph{both}
baselines (welded and imperfect), exactly as H2 requires --- comparing only
against the welded case would flatter the interphase.
\end{remark}

\begin{remark}[Sign and frame conventions]
The sign in front of the bracket in \eqref{eq:h0} is fixed by
\eqref{eq:thermal}: with the ordering $\jmp{f}=f^{(2)}-f^{(1)}$ a positive
compliance produces a negative flux jump for positive flux. Reversing the jump
ordering flips this sign, so the ordering must be declared once, as it is in
\eqref{eq:jump}.
\end{remark}

\section{Independent-equation count}
\label{sec:count}

Blueprint \S4 requires the count to be \emph{derived per configuration}. The
derivation proceeds in three steps: count unknown amplitudes, count independent
relations, check for gauge degeneracies.

\subsection{Step 1: admissible amplitudes}

From \texttt{theory/characteristic\_matrix.tex}, the characteristic determinant
has $8$ roots in $\xi$, forming $4$ pairs $(\xi,-\xi)$ by evenness. Of the $8$
roots, the number admissible on a given side is determined by the decay/outgoing
selection: for medium 1 ($z<0$) a mode is admissible if $\operatorname{Im}\xi<0$
(decays as $z\to-\infty$), and for medium 2 if $\operatorname{Im}\xi>0$. For real
$\xi$ the sign of $\operatorname{Re}\xi$ decides. Exactly $4$ roots satisfy each
condition, so:
\begin{center}
\begin{tabular}{lll}
\toprule
Region & Admissible amplitudes & Basis \\
\midrule
Medium 1 (reflected) & 4 & 4 outgoing roots of $M_1$ \\
Medium 2 (transmitted) & 4 & 4 outgoing roots of $M_2$ \\
Interphase, forward & 4 & 4 outgoing-looking roots in $0<z<\mathrm{h}$ \\
Interphase, backward & 4 & the same 4, traversed the other way \\
\bottomrule
\end{tabular}
\end{center}
Plus the incident amplitude, which is a fixed datum (normalisation), not an
unknown.

\subsection{Step 2: independent relations per face}

At one face there are eight conditions:
\begin{equation}
\begin{aligned}
&\text{mechanical: } &&\jmp{\s_{xz}}=0 \text{ or } \s_{xz}=K_x\jmp{u_x}+C_x\jmp{\dot u_x}, \\
& &&\jmp{\s_{zz}}=0 \text{ or } \s_{zz}=K_z\jmp{u_z}+C_z\jmp{\dot u_z}, \\
&\text{electrical: } &&\text{closure E-A or E-B of \S\ref{sec:elec} (2 relations)}, \\
&\text{thermal: } &&q_z^{(1)}=q_z^{(2)}=q_c \text{ with } q_c=-H_T\jmp{\bth} \text{ (2 relations)}.
\end{aligned}
\label{eq:eight}
\end{equation}
That is $8$ relations per face.

\subsection{Step 3: counts and gauge removal}

\begin{center}
\begin{tabular}{lccc}
\toprule
Configuration & Unknowns & Relations & Solvability \\
\midrule
Two half-spaces, zero-thickness contact & $4+4=8$ & $8$ (one face) & square \\
Interphase, $0<z<\mathrm{h}$ & $4+4+4+4=16$ & $2\times8=16$ & square \\
Welded limit, two half-spaces & $4+4=8$ & $8$ & square \\
\bottomrule
\end{tabular}
\end{center}

The counts are square in every configuration, so the systems are generically
determined. This is a necessary condition for well-posedness, \textbf{not a
sufficient one}, and the distinction is important: a square system can still be
rank-deficient, and a rank-deficient system can still have a full-rank operator
with wrong normalisation. Check I2 computes the rank of the assembled operator
for the welded piezothermoelastic interface at three angles and obtains $8$ in
every case.

\subsection{Gauge degeneracies, and why they matter}

Two degeneracies must be removed before the count is trusted.

\begin{remark}[Gauge 1: constant potential at $k=0$]
At $(\kappa,\xi)=(0,0)$ the null space of $M$ is one-dimensional and is the
constant-potential state $(0,0,1,0)$, which carries $D_z=0$ and therefore no
physics. Its presence means normal incidence must be taken as the limit
$\kappa\to0$, never evaluated at $\kappa=0$. \textsc{Derived} and \textsc{Tested}
(check C1).
\end{remark}

\begin{remark}[Gauge 2: the exactly decoupled limit is singular]
\label{rem:decoupled-singular}
If the piezoelectric, thermoelastic and pyroelectric couplings are set to
\emph{exactly} zero, the potential $\varphi$ and the temperature $\bth$ appear in
no equation other than their own. Four of the eight conditions \eqref{eq:eight}
then become identically zero rows --- the interface system is singular and has no
unique solution, even though the physics of the decoupled problem is perfectly
well posed. \textbf{The decoupled case must be reached as a limit}, with the
couplings scaled by a small parameter $\lambda\to0$, and the solution converges
as $O(\lambda^2)$ because the piezoelectric contribution to the longitudinal
modulus is $e^2/\varepsilon$.
\end{remark}

This second degeneracy was found the hard way: an attempt to validate the solver
against the elastic Fresnel result by setting $e=\beta=p=0$ produced a singular
$8\times8$ system at every parameter point. It is recorded as item U-3 in
\texttt{reports/UNRESOLVED\_ISSUES.md}.

\section{Summary}
\label{sec:summary}

\paragraph{Established.}
The contact laws \eqref{eq:spring}, \eqref{eq:elec-EA}/\eqref{eq:elec-EB} and
\eqref{eq:thermal}; the dissipation forms \eqref{eq:mech-nonneg},
\eqref{eq:therm-diss}, with non-negativity proved and units checked; the
reactive character of the capacitive gap; the first-order $\mathrm{h}\to0$
reduction \eqref{eq:h0}; and the equation counts $8$/$16$ with rank $8$ verified.

\paragraph{Not established.}
\begin{itemize}
  \item The overall factor convention ($\tfrac12$ or not) for the thermal
        dissipation, which must be unified with the mechanical convention before
        any numerical energy balance (U-5).
  \item The $O(\mathrm{h}^2)$ term and the non-diagonal form of \eqref{eq:h0}
        (U-6).
  \item Whether any real interface has a conductive electrical path as opposed to
        a purely capacitive one; Phase 1 asserts neither.
  \item No cross-coupling between the mechanical, electrical and thermal contact
        channels is introduced anywhere. Blueprint \S4 forbids a full $4\times4$
        off-diagonal contact block without an interface potential and a
        dissipation law; none is available, so the block is kept diagonal.
\end{itemize}

\end{document}
